% Source: https://artofproblemsolving.com/community/c3h1355615_trigonometry
% Stable identity: aops-topic-1355615-q01
On the unit circle, $\cos108^\circ$ is the $x$-coordinate of the point $P$ at angle $108^\circ$; it is negative.

\begin{asy}
usepackage("amsmath");
size(4.5cm);
defaultpen(fontsize(10pt));
real c=(1-sqrt(5))/4;
pair O=(0,0),P=dir(108),Q=(c,0);
draw(circle(O,1),gray);
draw((-1.25,0)--(1.25,0)); draw((0,-1.2)--(0,1.25));
draw(O--P,blue); draw(P--Q,dashed);
draw(arc(O,0.25,0,108),red); label("$108^\circ$",0.42*dir(54),red);
dot(P);dot(Q);
label("$P$",P,NW); label("$O$",O,SE); label("$\cos108^\circ$",(-0.42,0),S);
label("$1$",O+0.72*(P-O),E,blue);
label("\mbox{}",(min(currentpicture,user=true)+max(currentpicture,user=true))/2,black); // keeps the text after the figure black when LaTeX typesets the labels inline
\end{asy}

Let $\theta=36^\circ$, so $3\theta=180^\circ-2\theta$ and $\cos3\theta=-\cos2\theta$. With $c=\cos\theta$, the formulas $\cos3\theta=4c^3-3c$ and $\cos2\theta=2c^2-1$ give
\[
4c^3-3c=-(2c^2-1)\iff4c^3+2c^2-3c-1=0\iff(c+1)\left(4c^2-2c-1\right)=0 .
\]
Since $c=\cos36^\circ>0$, $c=\dfrac{2+\sqrt{20}}{8}=\dfrac{1+\sqrt5}{4}$. Then
\[
\cos72^\circ=2c^2-1=2\cdot\dfrac{6+2\sqrt5}{16}-1=\dfrac{\sqrt5-1}{4},
\]
and $\cos108^\circ=\cos(180^\circ-72^\circ)=-\cos72^\circ=\boxed{\dfrac{1-\sqrt5}{4}}\approx-0.309$.

The special values of $30^\circ,45^\circ,60^\circ$ alone are not enough: sums and differences of those angles give only multiples of $15^\circ$, and $108^\circ$ is not one. A polynomial equation such as the one above is needed.
